\section*{第 \ref{ch:b3:genetic-engineering} 章——基因工程与生物技术}

\begin{solution}{exo:b3:genetic-engineering:1}
一个复制起点，使宿主能复制\omterm{def:b3:genetic-engineering:tools}{质粒}，
并将其传给子细胞；一个\omterm{def:b3:genetic-engineering:tools}{选择标记}，通常为抗生素抗性，
使只有携带\omterm{def:b3:genetic-engineering:tools}{质粒}的细胞得以生长；以及一个唯一的限制位点（或一簇位点），
插入片段在此连接，
且不在\omterm{def:b3:genetic-engineering:tools}{质粒}其他位置切割。
\end{solution}

\begin{solution}{exo:b3:genetic-engineering:2}
在随机 DNA 中，给定的六碱\omterm{def:b3:bioinformatics:motif}{基序}列每
$4^{6} = 4096$ 个碱基对出现一次，故约 $4.6\times 10^{6}/4096 \approx 1100$ 个片段，平均长
\qty{4.1}{kb}。真实基因组并非随机——碱基组成、
密码子使用偏好以及对某些序列的回避都会改变这一数目（
\emph{E.~coli} 中 EcoRI 位点的实际数目为数百个）。
\end{solution}

\begin{solution}{exo:b3:genetic-engineering:3}
变性温度为 \qty{95}{\celsius}，使双链分开；退火
温度为 \qtyrange{50}{65}{\celsius}，使\omterm{def:b3:genetic-engineering:pcr}{引物}与其结合位点配对；
延伸温度为 \qty{72}{\celsius}，使聚合酶从每条
\omterm{def:b3:genetic-engineering:pcr}{引物}开始复制。普通聚合酶在 \qty{95}{\celsius} 下
会在第一个循环中失活，需反复补加三十次；Taq
酶则能耐受全部循环。
\end{solution}

\begin{solution}{exo:b3:genetic-engineering:4}
一条向导 RNA，其二十个碱基与靶序列匹配，另有 \omterm{def:b3:genetic-engineering:crispr}{PAM}（NGG）
紧接靶序列 3$'$ 端、位于非靶链上。\omterm{def:b3:dna-repair:lesions}{双链断裂}后，
可发生末端连接，留下小的插入或
缺失，通常会破坏基因的阅读框；或发生\omterm{def:b3:dna-repair:dsb}{同源重组}，
借助携带所需序列的模板进行，
将其写入。
\end{solution}

\begin{solution}{exo:b3:genetic-engineering:5}
$50\times 1.9^{35} = 50\times 5.7\times 10^{9} \approx 3\times 10^{11}$
个拷贝。$\Delta C_{T} = 7.3$；比值 $1.95^{7.3} = e^{7.3\ln 1.95}
\approx 130$：第一个样本的模板约多 $130$ 倍。
\end{solution}

\begin{solution}{exo:b3:genetic-engineering:6}
细菌无法剪接：带有内含子的基因组拷贝
转录后将得到无用的信使，因此改用无内含子的 cDNA，
由成熟 mRNA 反转录而来。其他障碍包括：人源启动子
不能被细菌聚合酶识别（需提供细菌启动子及
核糖体结合位点）；人源密码子偏好不同（需按
宿主偏好的密码子合成该基因）；蛋白质可能无法折叠或
发生聚集（降低温度、与可溶性伴侣融合，或
从包涵体中复性）；缺乏糖基化（可改用酵母或哺乳动物
细胞，若有必要）；二硫键需在氧化性周质中形成。
\end{solution}

\begin{solution}{exo:b3:genetic-engineering:7}
将靶向修饰的干细胞注入宿主囊胚，并
与其自身细胞混合，因此小鼠由两种基因型构成——即
嵌合体——只有当其部分生殖细胞来自靶向细胞
时，它才会传递该等位基因。将嵌合体与
野生型小鼠交配，可得杂合子（约一半配子
携带该等位基因，前提是生殖系一半被靶向修饰）；杂合子相互交配，可得四分之一
纯合\omterm{def:b3:genetic-engineering:transgenic}{敲除}小鼠。
\end{solution}

\begin{solution}{exo:b3:genetic-engineering:8}
精确估计：$6.4\times 10^{9}\times 4^{-20}\times (1/16) \approx 4\times
10^{-4}$ 个位点（\omterm{def:b3:genetic-engineering:crispr}{PAM} 中的 GG 使其再降低 $16$ 倍）——实际
上几乎没有。两处错配以内，共有 $1 + 60 + 1710 = 1771$
个序列，$1771\times 4\times 10^{-4} \approx 0.7$ 个偶然位点。真实
基因组并非随机，而向导序列会直接与之核对。
\end{solution}

\begin{solution}{exo:b3:genetic-engineering:9}
$p_{1} = 0.05 + 0.8\times 0.05\times 0.95 = 0.088$; $p_{2} = 0.152$;
$p_{3} = 0.255$；随后为 $0.41$、$0.60$：约五代后
超过 $0.5$，对照估计值 $\ln(1/(2\times 0.05))/\ln 1.8 \approx 4$。
\end{solution}

\begin{solution}{exo:b3:genetic-engineering:10}
末端连接在各类细胞中都起作用，包括处于静止状态的干细胞，
而\omterm{def:b3:dna-repair:dsb}{同源重组}是 S/G2 期的过程，在其中效率低下；它
无需递送模板，其产物——任何对
增强子的破坏——均可奏效，而矫正则需要精确的
序列，且只能在少数细胞中实现。另外，同一编辑可
适用于每一种 $\beta$-珠蛋白突变，无论是镰状细胞型还是地中海贫血型。缺点是它
去除了一个调控元件（及其所承担的其他作用），它
使镰状等位基因原样保留，插入缺失的混合情况
不受控制。
\end{solution}

\begin{solution}{exo:b3:genetic-engineering:11}
设纯合子适合度为 $1 - s$：$p' = \bigl[p^{2}(1-s) + p(1-p)(1+c)
\bigr]/(1 - sp^{2})$。在等位基因稀有时纯合子可忽略，且
驱动等位基因按 $(1+c)p$ 增长，无论 $s$ 为何值；驱动是否会
达到固定，由 $p = 1$ 附近的情况决定，此处频率为 $q$ 的
稀有野生型等位基因回归为 $q' \approx q(1-c)/(1-s)$：若 $c > s$ 则固定，
否则为一个中间平衡。抗性：每次失败的转换
（改用末端连接而非重组）会留下一个向导不再识别的
等位基因；若此类等位基因具有适合度——即位于非必需区域的框内
插入缺失——则它们不付代价，而驱动却付出代价，因此选择
有利于它们，它们将取代驱动。补救办法是
靶向插入缺失会致死的必需序列，并同时使用多个
向导。
\end{solution}

\begin{solution}{exo:b3:genetic-engineering:12}
注入受精卵的 \omterm{def:b3:genetic-engineering:crispr}{Cas9} 可能在首次分裂之后才发挥作用，因此
不同细胞接受不同的修复——即镶嵌现象；每次断裂的修复
由细胞选择，末端连接产生不可预测的
插入缺失，以及切口处的大片段缺失或杂合性丢失；
\omterm{def:b3:dna-repair:dsb}{同源重组}的结果仅占少数；\omterm{prop:b3:genetic-engineering:editors}{脱靶}断裂
发生在无法全部预测的位点上；而胚胎只能
通过对少量活检细胞测序来检测，这无法代表
其余细胞。所需证据：能让每个细胞都以相同方式编辑的方法
（在首次 S 期之前），效率接近 \qty{100}{\%}，
对编辑后动物胚胎的每个细胞进行全基因组测序，显示
无非预期的改变，并对编辑动物进行长期随访，
跨代观察——这些证据目前均不存在。
\end{solution}

\begin{solution}{pb:b3:genetic-engineering:1}
\textbf{1.} 每 $4096$ bp 出现一次。$P(\text{无酶切位点}) = (1 - 1/4096)^{1500}
= e^{-0.37} = 0.69$.
\textbf{2.} 选用基因中无酶切位点的\omterm{def:b3:genetic-engineering:tools}{限制酶}；或
用 PCR 扩增该基因，所用\omterm{def:b3:genetic-engineering:pcr}{引物}在其 5$'$
端带有新的酶切位点（或使用平末端，或不依赖连接的克隆）。
\textbf{3.} \qty{0.1}{\micro g}\,$\times 2\times 10^{7} = 2\times
10^{6}$ 个菌落；\qty{10}{\%}，即 $2\times 10^{5}$ 个携带插入片段。
\textbf{4.} 白色菌落：\qty{10}{\%}。插入片段中一半
方向正确，故 $P(\text{ 在 } n \text{ 个白色菌落中均无 }) = 0.5^{n} \le 0.01$
得 $n \ge 6.6$，故取七个。
\textbf{5.} $32$ 次倍增，$2^{32} = 4.3\times 10^{9}$ 个细胞，$2.1\times
10^{11}$ 个\omterm{def:b3:genetic-engineering:tools}{质粒}；每个\omterm{def:b3:genetic-engineering:tools}{质粒}的质量为 $5500\times 650 = 3.6\times 10^{6}$ Da $=
5.9\times 10^{-18}$ 克：\qty{1.3}{\micro g} \omterm{def:b3:genetic-engineering:tools}{质粒}。
\textbf{6.} 克隆只需复制，由\omterm{def:b3:genetic-engineering:tools}{质粒}的
复制起点提供；人类启动子不能被细菌 RNA 聚合酶识别，因此
要使其表达，须在
（无内含子的）编码序列前加入细菌启动子和核糖体结合位点。
\textbf{7.} $N_{0} = N_{T}/2^{C_{T}}$: $10^{12}/2^{28} \approx 3700$
个拷贝；当 $C_{T} = 35$ 时，$10^{12}/2^{35} \approx 29$ 个拷贝。
\textbf{8.} $2^{n} = 10^{12}$：$n = 39.9$，即 40 个循环。
\textbf{9.} $10^{12}/2^{38} \approx 3.6$ 个拷贝。超过约 $40$
个循环后，单个污染分子或\omterm{def:b3:genetic-engineering:pcr}{引物}假象即可
达到阈值，而拷贝数不足一个的样本毫无意义。
\textbf{10.} 仅一个分子在循环 $40$ 时即越过阈值，判为阳性。防止此类情况，
应将配制与处理产物分设房间和设备，
采用单向流程，每次运行均设阴性对照，
并以酶法破坏残留的扩增子。
\textbf{11.} $3700/0.4 \approx 9000$ 个 RNA 拷贝。
\textbf{12.} 报告的倍数为 $2^{3.3} \approx 10$ 倍；真实倍数为 $1.9^{3.3} =
e^{3.3\ln 1.9} \approx 8.3$ 倍。
\textbf{13.} $40\times \qty{35}{\micro g} = \qty{1.4}{mg}$ 每天，
\qty{0.51}{g} 每年；$\times 10^{7}$：\qty{5100}{kg} 每年。
\textbf{14.} $5.1\times 10^{6}/5800 \approx 880$ 摩尔；$5.3\times 10^{26}$
个分子。
\textbf{15.} $5.1\times 10^{6}$ 克 $/ (20$ 克/升$) = 2.6\times 10^{5}$ 升
$= \qty{260}{m^3}$——为池子容量的十分之一。
\textbf{16.} $5100\times 8000 = 4.1\times 10^{7}$ 千克胰腺，按
每个 \qty{0.1}{kg} 计，每年需 $4\times 10^{8}$ 头猪。
\textbf{17.} 重组激素与人体激素完全相同，故它
不会诱导抗体产生，也不引起过敏反应；它不携带
动物病原体（病毒、朊病毒）；其供应不依赖屠宰；
且序列可以得到改进。
\textbf{18.} 与受体接触的残基未变，因此
结合和信号传递与胰岛素相同；改变的残基位于
胰岛素分子表面，即胰岛素分子缔合成二聚体和
六聚体的表面，改变它们会改变注射储库
解离为可吸收单体的速度。
\textbf{19.} $6.4\times 10^{9}\times 4^{-20}/16 \approx 3.6\times
10^{-4}$：无精确的偶然匹配。
\textbf{20.} $1 + 60 + 1710 + 27\times 1140 = \num{32551}$ 条序列；
$\num{32551}\times 3.6\times 10^{-4} \approx 12$ 个偶然的\omterm{prop:b3:genetic-engineering:editors}{脱靶}
位点。
\textbf{21.} $10^{7}\times 10^{-4} = 1000$ 个细胞。干细胞
在患者一生中存活并分裂；某个细胞若丢失了
抑癌基因，或获得了易位，即可建立一个将
发展为白血病的克隆。
\textbf{22.} $p_{1} = 0.0019$, $p_{2} = 0.0036$, $p_{3} = 0.0069$;
每代以 $1.9$ 倍几何增长，由此得 $\ln 500/\ln 1.9
\approx 10$；迭代递推式，频率超过 $0.5$ 的时刻
为第 $11$ 代（$0.45 \to 0.68$）。
\textbf{23.} 约一年。纯合子中 \qty{20}{\%} 的适合度代价小于
$c = 0.9$，故驱动仍会固定，速度较慢，同时
种群平均适合度下降——这正是
抑制型驱动的目的；若代价高于 $c$，驱动将停滞于中间
频率。
\textbf{24.} $10^{6}\times 0.1 = 10^{5}$ 个抗性等位基因，于一
代中产生。抗性等位基因不可切割，若适合度良好，则在驱动有代价时
无代价，因而受到选择并扩散，驱动随之被消除——因此
故驱动应靶向末端连接产物致死的位点，
并同时使用多条向导 RNA。
\textbf{25.} 首份样本中约 $3700$ 个拷贝；每年约需 \qty{260}{m^3}
的发酵罐以供全球胰岛素生产；驱动约需 $11$ 代，
一个蚊子年，即驱动所需的时间。
\end{solution}
